\section{Conclusion and Outlook}

In this report, we studied Signorini's problem for a linearly elastic body in unilateral, frictionless contact with a rigid foundation. Starting from the geometrical and physical setup, we introduced the relevant Sobolev and trace theoretic framework and formulated the problem first in terms of admissible displacements and elastic energy. From there, we derived and compared several equivalent formulations, namely the {\em energy formulation}, the {\em weak formulation} as a variational inequality, stronger formulations involving normal traces and contact conditions, and finally a {\em mixed formulation} with a Lagrange multiplier representing the normal contact force. In this way, the report aimed to make precise how the different viewpoints on Signorini's problem are related and how the contact conditions can be recovered in increasingly explicit form.

At the same time, the present work only treats the frictionless case. A natural extension is to include tangential contact effects, i.e.\ friction, on the contact boundary. In the book \cite{book}, this is pursued first through the {\em Tresca friction model} and afterwards through {\em Coulomb friction}. From the analytical point of view, this leads to significantly more difficult problems. While the frictionless Signorini problem still fits into the relatively well-structured setting of a variational inequality of the first kind, friction introduces additional nonlinearities and, in the Coulomb case, even a quasi-variational structure. Thus, the frictionless setting considered here may be viewed as the basic model from which more realistic, but also more delicate, contact models can be developed.

Another natural continuation concerns numerical approximation. Since the Signorini problem already admits a weak and a mixed formulation, it fits naturally into the framework of {\em finite element methods}. In particular, the variational inequality formulation provides the starting point for conforming finite element discretizations, while the mixed formulation is well suited for methods involving {\em Lagrange multipliers} for the contact force. Beyond this, the book also discusses {\em Nitsche-type methods} and {\em augmented Lagrangian techniques} as alternative approaches.